\chapter{Theoretical Background}
This chapter provides the technical background that is required to understand the concepts of this thesis properly.
There are many terms used in the vision--language literature that might seem similar to each other, but mean different things.
For example, answer accuracy, explanation quality, and visual grounding are related concepts, but they are measured differently.
Also, concepts like contrastive alignment, generative fine-tuning, and explanation-aware training are all about connecting the visual signal to the language model,
but they do it through different training signals and have different implications for the model's behaviour.
The main goal of this chapter is to provide enough background on these concepts so that the reader understands
what they mean, and how and why they are used in the experiments of this thesis.

\section{Vision--language models and visual question answering}
A vision--language model is a model that can process both visual and text-based inputs and produce text-based outputs.
In this thesis, the visual input can be a natural image, a chart, a scanned document, or a science-style diagram.
The text input is usually a question to be answered based on the provided image, and in some datasets it also includes answer choices.
The model is expected to generate an answer to the given question, and in some cases it is also expected to produce a rationale before answering the question.

The task throughout this thesis is visual question answering.
The structure is simple: given an image and a question as input, the model returns an answer as output.
However, the type of reasoning required to come up with the answer can be different across datasets.
In a natural-image question, the answer usually depends on recognising an object or a simple relationship between objects.
In the case of a chart question, the answer might depend on reading an axis label, comparing two bars, or estimating a value.
In a document question, the answer might be a short text region hidden somewhere inside a scanned page.
In a science or commonsense question, the model might need to combine image information with answer choices and learned background knowledge.

This variation is one of the reasons why this thesis uses multiple datasets across different domains instead of focusing on one dataset only.
A model might be good at one type of visual reasoning but perform poorly on another type of visual reasoning.
For example, a model that excels at recognising objects in natural images might still fail to read small text in a document or to compare two bars in a chart.
That is why this thesis treats visual question answering as a family of related tasks rather than one single problem.

A second important point is that the answer format is different across datasets.
In ScienceQA and A-OKVQA datasets the answers are multiple-choice, where the model does not automatically generate the answer,
but instead assigns a score to each answer choice and the highest-scoring choice is then selected as the final answer.
Other datasets, such as DocVQA, ChartQA, and VQAv2 provide open-ended or short-answer tasks, where the model produces the answer directly.
That is why the evaluation cannot use a single universal metric for every dataset and instead uses the metrics
that are native to each dataset.

The wider VQA literature shows why these task differences matter.
VQA and OK-VQA focus on natural-image and knowledge-heavy question answering,
TextVQA and OCR-VQA add text in the image,
GQA and CLEVR focus more on compositional reasoning,
and DVQA, FigureQA, PlotQA, and InfographicVQA move the same problem toward charts, plots, and infographic-style documents~\cite{agrawal2017vqa,marino2019okvqa,singh2019textvqa,mishra2019ocrvqa,hudson2019gqa,johnson2017clevr,kafle2018dvqa,kahou2018figureqa,kim2020plotqa,mathew2022infographicvqa}.


\section{Cross-modal alignment}
Cross-modal alignment means connecting signals from different modalities, such as images and language,
in a way that those modalities can interact with each other for the model to perform a certain task.
In vision--language models, the visual signal must provide useful context for the language side.
If this connection is weak, the model might still produce a convincing answer, but that answer might be coming from
the learned language patterns rather than from the visual information provided by the image.

The word alignment can mean different things in different contexts.
At the representation level, it means that the model learns that in the embedding space, the provided image and the matching text should be close to each other, while non-matching ones should be further apart.
At the generation level, it means that the model learns to use the information provided by the image when generating the output for a given task.
At the reasoning level, it means that the model learns to produce explanations that connect the question, visual input, and the final answer in a helpful way.

Earlier alignment models such as CLIP and ALIGN make the representation-level idea very clear:
matching images and text are pulled together in the embedding space,
while BLIP and Flamingo push the same alignment idea toward generation and few-shot vision--language interaction~\cite{radford2021clip,jia2021align,li2022blip,alayrac2022flamingo}.
Grounding datasets such as Visual Genome, Flickr30k Entities, ReferItGame, and RefCOCO also show the stricter version of the problem,
where the model has to connect language not only to the whole image, but also to the right region inside the image~\cite{krishna2017visualgenome,plummer2015flickr30k,kazemzadeh2014referitgame,yu2016modelingvisual}.

A model might internally have a strong text--image representation and still fail to generate a correct answer.
It might generate a correct answer without producing a useful explanation.
It might also produce a good-looking explanation that is not actually based on the right visual evidence.
That is why this thesis separates the training objective from the evaluation axis.
The training objective tells us what signal to feed the model during fine-tuning,
while the evaluation axis tells us which behaviour actually changed after fine-tuning. 

During comparison, the model design, the data handling, and the compute budget are kept constant, so that the outcome is attributed to the training objective rather than to other factors.
If one method uses a larger model, more training data, or a different prompt format, then it would be difficult to say whether the outcome is the result of the training objective or any of those other factors.
This does not make the study perfect, but it makes the comparison more focused on the effect of the training objective itself.


\section{Generative supervision}
Generative supervision is a setup in which the model is trained to produce a target token sequence.
In the answer-only case, that target sequence is comprised of the answer tokens only.
In a rationale-generative case, the target sequence contains not only the answer tokens, but also a reasoning sentence or rationale before the answer.
The model is fed an image and a question as input, and the loss is computed on the target tokens.

Since instruction-tuned VLMs already work by reading a prompt and generating text,
training them with answer/rationale text targets fits well with the way these models are used.
The model can already generate text, so fine-tuning can adapt it to produce the expected output format for a specific task.
For example, the model can learn that in ScienceQA, it should select an answer from the provided answer choices,
or that in DocVQA, it should generate a short text as the answer.
Hence, generative fine-tuning is a good baseline for many visual question answering tasks.

However, pure generative supervision has some limitations when it comes to grounding.
A model might learn to generate the answer without actually using the provided visual information correctly.
It might also just memorise common patterns in the dataset.
For example, in multiple-choice datasets, some questions could to some extent be answered by looking at the question and the answer choices alone,
without actually looking at the image.
In documents and charts, the model could memorise that the answer is often in some specific position or format.
This does not mean that generative supervision is not a good option, but it does mean that in the evaluation phase, we need to check more than just the final answer accuracy
to make sure that the model is actually using the visual signal in a meaningful way.

In this thesis, we use generative supervision in two ways.
First, we use it as a direct training objective for the BLIP-2 and Qwen2-VL runs that are not explanation-aware.
Second, we use it as a control for the explanation-aware training.
This is quite important, because we do not want to compare the explanation-aware supervision only against a frozen zero-shot model,
but also against pure generative supervision that was fine-tuned to produce the expected answer or rationale-plus-answer text without
the same span-weighting strategy. 


\section{Contrastive supervision}
Contrastive supervision is a way to train the model to bring the image and text representations closer to each other when they match,
and to push them apart when they do not match.
In this thesis, we use contrastive supervision as an additional signal on top of generative fine-tuning, not as a replacement for it.
We still train the model to generate the answer, but in addition to that, we also give it a contrastive signal over image--answer representations.
We use BLIP-2 for this setup, because it has a Q-Former that uses a fixed number of learnable query tokens to pool visual information from the image encoder outputs, which can then be used to build an image representation.
The basic idea is that the image representation should be closer to the correct answer representation than to other answers in the same batch.

This contrastive-enhanced setup is different from ordinary text generation.
The generative part asks the model to generate the next target token.
The contrastive part asks the model to organise the image and answer representations in such a way that the correct image--answer pair is easier to retrieve than mismatched pairs.

In principle, this can encourage the model to learn stronger image--answer alignment.
In practice, however, the effect of the contrastive signal depends on many factors, such as the model architecture, the way the image representation is extracted,
the design of the projection head for contrastive learning, the batch size used during training, and the nature of the data.
Hence, we do not assume that the contrastive branch is useful just because it adds an additional loss term.
Instead, we test the contrastive-enhanced setup directly against a BLIP-2 generative fine-tuning control, to check whether it actually leads to better results.

One of the important things to check is the retrieval diagnostic.
If the contrastive supervision is doing what we expect, the correct answer should be easier to retrieve from the image representation than mismatched answers.
A simple way to check that is to compute Recall@1, Recall@5, Recall@10, and mean reciprocal rank.
These metrics do not directly measure the final answer accuracy, but they give us a way to see if the contrastive objective is actually adding any value in terms of representation alignment.

A modest contrastive result does not mean that contrastive supervision is useless in general.
It only means that under our specific setup, with the chosen BLIP-2 configuration and the capped ScienceQA run, the contrastive-enhanced fine-tuning
does not lead to a better result than the generative fine-tuning control.
This is important to note, because the thesis is about the controlled effect of the training objective under the available budget, not about dismissing an entire class of methods.


\section{Rationales and explanation-aware learning}
A rationale is a text generated before the final answer that explains how to connect the question and the visual input to come up with the answer.
Not all the datasets provide gold rationales.
Hence, for the rationale supervision, we use only the datasets that provide gold rationales, which are ScienceQA and A-OKVQA.
With those runs, we can check whether rationale supervision improves answer accuracy and explanation quality.

Explanation-aware supervision is very similar to rationale supervision with the only difference in the training loss.
In rationale supervision, we train the model to generate a rationale and an answer in one sequence, and the loss is computed over the entire sequence.
In explanation-aware supervision, we split the target into two spans: the explanation span and the answer span.
The training loss is then computed separately for each span, and a coefficient $\alpha$ is used to weight the answer loss compared with the explanation loss.
This allows us to test whether answer supervision and rationale supervision help each other or compete with each other.

There is an ablation sweep done over $\alpha$ to check how the balance between answer and explanation supervision affects the results.
When $\alpha$ is low, the model receives more signal from the explanation span,
which encourages it to produce better rationales, but it might not focus enough on getting the answer right.
On the other hand, when $\alpha$ is high, the model receives more signal from the answer span,
which encourages it to produce better answers, but it might not focus enough on producing good rationales.
The third setting is a length-aware control,
where $\alpha$ is set based on the number of tokens in the explanation and answer spans.
This control is useful because fixed $\alpha$ values might give a misleading impression of what explanation-aware training is doing.

A key limitation is that only datasets with gold rationales can directly support rationale-supervised training.
ScienceQA and A-OKVQA provide this supervision in the completed experiments.
ChartQA, DocVQA, and VQAv2 do not provide the same type of gold rationale in this setup.
For those datasets, the thesis uses answer-only fallback training and transfer testing instead of treating them as explanation-aware datasets.
This is why the interpretation of results changes by domain.

Rationale quality is also not the same as faithfulness.
A generated explanation can overlap with the gold rationale and still fail to prove that the model actually used the image.
For example, the model might learn a common explanation pattern from training data.
It might also generate a reasonable explanation after choosing the answer for some other reason.
This is why rationale metrics are reported beside, not instead of, evidence-sensitivity analysis.

\section{Parameter-efficient fine-tuning}
Full fine-tuning updates all or most of the parameters in a model.
For large vision--language models, this can be expensive in memory, compute, and time.
Parameter-efficient fine-tuning updates only a small number of additional parameters while leaving most of the base model frozen or quantised.
This makes experiments possible on a much smaller budget.

LoRA is one common parameter-efficient method.
Instead of updating a full weight matrix directly, LoRA adds low-rank trainable matrices that modify the behaviour of selected layers.
QLoRA extends the idea by loading the base model in quantised precision and training low-rank adapters on top of it.
This makes the memory requirement lower while still allowing the model to adapt to the task.
In this thesis, Qwen2-VL runs use QLoRA so that both the broad sweep and the 7B scale check can fit the available environment.

LoRA and QLoRA are part of a larger family of parameter-efficient methods.
Adapters, prefix tuning, prompt tuning, P-tuning v2, BitFit, and vision--language adapters all follow the same practical direction:
change a small trainable part of the model instead of updating the whole backbone~\cite{houlsby2019adapters,li2021prefix,lester2021prompt,liu2022ptuningv2,zaken2022bitfit,sung2022vladapter}.

Parameter-efficient fine-tuning is not only a practical trick.
It also shapes the research question.
The thesis asks what can be learned under a resource-constrained setup, not what could be achieved with unlimited compute and full-model fine-tuning.
That is why the results should be read as fixed-budget comparisons.
A method that works well under these constraints is useful for students, small labs, and early-stage research projects.
A method that only works after very large compute investment would answer a different question.

Using QLoRA also makes the scale comparison more controlled.
The 2B and 7B Qwen2-VL runs are not compared by changing the whole training strategy.
Both use the same explanation-aware setting and the same parameter-efficient adaptation style.
This makes the observed difference more about backbone scale than about changing the fine-tuning method.
The comparison is still not a full efficiency benchmark because exact runtime and peak memory are not measured as primary results.

\section{Evidence sensitivity and masking}
Evidence sensitivity asks whether the model's answer changes when relevant visual information is removed or altered.
In this thesis, the diagnostic is based on masking regions of the image.
The model is scored on the full image, then scored again after a region is hidden.
If the score for the correct answer drops after masking, the model was sensitive to that region.

This is related to the broader explainability literature.
Methods such as Grad-CAM, LIME, Integrated Gradients, and Transformer attribution also try to connect model behaviour back to input evidence,
but this thesis uses a simpler score-change masking diagnostic instead of claiming a full causal explanation~\cite{selvaraju2017gradcam,ribeiro2016lime,sundararajan2017axiomatic,chefer2021transformer}.

This is a useful diagnostic, but it must be interpreted carefully.
A high drift value means the answer score depends on the masked region.
It does not automatically prove that the generated rationale is faithful.
A model may rely on a visual region without describing it correctly.
It may also produce a good-looking rationale while the answer score barely changes when the image is masked.
For that reason, masking drift is treated as evidence sensitivity rather than a full proof of faithful explanation generation.

Masking is also affected by the domain.
In documents and charts, evidence can be local: a number in a table, a word in a scanned form, an axis label, or a bar height.
Masking such a region can strongly change the answer score.
In science and commonsense images, the evidence may be more diffuse, or the answer may be partly recoverable from the question and answer choices.
In that case, masking one coarse region may produce a smaller drift even when the image still matters.

The thesis therefore uses masking as one part of the evaluation, not as the only faithfulness measure.
It is strongest when read together with answer accuracy, rationale metrics, and qualitative examples.
The qualitative masking figure is included to show what kind of visual perturbation was applied and how the scores changed.
It is not meant to replace a human annotation study of faithful reasoning.

\section{Metrics used in the thesis}
The datasets in this thesis use different native metrics, so the results are not read as one universal score.
ScienceQA and A-OKVQA use multiple-choice accuracy in the main experiments.
ChartQA uses relaxed accuracy because numerical answers can have small acceptable differences.
DocVQA uses average normalised Levenshtein similarity, which gives partial credit for short text answers that are close to the target.
The VQAv2 subset uses VQA consensus accuracy where annotator answers are available.

The term headline metric refers to the main metric for a dataset.
It does not mean that a score on one dataset can be directly compared with a score on another dataset.
For example, 0.80 multiple-choice accuracy on ScienceQA and 0.80 ANLS on DocVQA do not mean the same thing.
They are both strong within their own benchmark context, but they measure different answer formats.
This is why the thesis compares models mostly within a dataset or within a clearly stated transfer column.

BLEU and ROUGE-L are used only when gold rationales are available.
They measure overlap between generated text and reference rationale text.
They are useful because they provide a repeatable automatic check, but they are limited.
A correct explanation can use different wording and receive a lower overlap score.
A high-overlap explanation can still fail to prove visual grounding.
For this reason, the thesis describes them as rationale quality metrics, not as complete faithfulness metrics.

Other text-quality metrics such as METEOR, BERTScore, and BLEURT are also common in generation evaluation,
especially when exact lexical overlap is too narrow.
They are not used as primary results here, because the goal is to keep the rationale-quality analysis lightweight and repeatable under the same fixed pipeline~\cite{banerjee2005meteor,zhang2020bertscore,sellam2020bleurt}.

The uncertainty intervals are also used cautiously.
Most completed evaluation sets are capped, and the main configurations use one seed.
Small differences such as 0.790 versus 0.795 should therefore be read as observed single-run differences.
They are not treated as conclusive evidence that one method is definitely better.
This is especially important in the BLIP-2 comparison and the A-OKVQA explanation-aware comparison, where the numerical gaps are small.

\section{Prompt families and model-facing text}
The prompt family is part of the experiment, not a small implementation detail.
If two models are trained or evaluated with different prompt styles, then a difference in performance may come from the prompt rather than from the training objective.
For that reason, this thesis separates answer-only prompts from rationale prompts and keeps this choice visible in the experiment ledger.
An answer-only run receives a question and is trained to produce only the final answer.
A rationale-generative or explanation-aware run receives the same visual question, but the target text contains a reasoning part before the final answer.

This distinction matters most for datasets with gold rationales.
ScienceQA and A-OKVQA can support rationale-supervised training because the dataset already contains explanations.
For those datasets, the model can be asked to produce a short reasoning trace and then the answer.
For ChartQA, DocVQA, and VQAv2, the completed setup does not have gold rationale labels.
Using rationale targets there would require an extra annotation or teacher-generation step, which would introduce a different source of supervision.
Therefore, those datasets are treated as answer-only in-domain and transfer settings.

The model-facing text is not exactly the same as the human-readable prompt written in the paper.
The processor inserts special tokens, image placeholders, and model-specific chat formatting.
This means that the actual tensor seen by the model is longer and more structured than the plain question shown in the dataset table.
The pipeline therefore saves sample traces that show the raw example, the rendered prompt, the target, the decoded token sequence, and the supervised spans.
Those traces are included as artifacts because they make the training process inspectable instead of hidden inside the collator.

This is also why the thesis avoids describing a run only as ``fine-tuned''.
Fine-tuning can mean different things depending on the prompt and target.
An answer-only adapter is trained on a different target than a rationale-generative adapter.
An explanation-aware adapter uses the same rationale-plus-answer target family, but it applies different loss weights to the answer and explanation spans.
So, when results are compared, the comparison has to say which prompt family and which target family were used.

\section{Collation and supervised spans}
Collation is the step where raw examples become tensors.
For this thesis, the raw example contains an image, a question, optional answer choices, a gold answer, and sometimes a gold rationale.
The collator turns that example into model inputs such as token IDs, attention masks, image tensors, and labels.
The labels tell the model which tokens should receive loss.
Prompt tokens are usually masked out, because the model should not be trained to reproduce the question.
Target tokens remain supervised, because they are the text the model is supposed to learn to generate.

This step is easy to underestimate.
In a normal text-only model, one might count the prompt length using a tokenizer and then start the labels after that point.
In a multimodal model, the processor may add image tokens before or inside the text representation.
If the label mask is shifted by even a few tokens, the model can be trained on the wrong span without an obvious runtime error.
The pipeline avoids this by measuring the prompt and full sequence through the same multimodal processor.
That makes the supervised answer and rationale spans line up with the sequence that is actually given to the model.

The span information is used by more than one objective.
For ordinary answer-only training, the supervised span is the answer target.
For rationale-generative training, the supervised span covers the rationale and answer target together.
For explanation-aware training, the rationale and answer spans are tracked separately so that the loss can weight them with alpha.
For the BLIP-2 contrastive-enhanced run, only the answer span is pooled into the contrastive text representation, while the rationale span still receives generative loss where it exists.
This keeps the contrastive signal focused on matching the image to the answer rather than to a longer explanation paragraph.

The saved sample traces are therefore not just debugging output.
They are part of the research record.
They show whether the answer-only run really excludes the rationale from the target.
They show whether the explanation-aware run actually has separate explanation and answer spans.
They also show whether decoded tokens match the intended prompt and target format.
Without this audit, a result could look mathematically valid while being based on the wrong supervision.

\section{Transfer evaluation}
Transfer evaluation asks what happens when an adapter trained in one source setting is evaluated on another target dataset or prompt family.
This is different from in-domain adaptation.
In-domain adaptation asks whether a model improves on the same dataset family it was trained on.
Transfer asks whether the learned behaviour carries over when the visual format, answer style, or prompt family changes.
For this thesis, that distinction is important because the domains are deliberately different.
Science questions, commonsense images, charts, documents, and general natural images do not stress the model in the same way.

The transfer matrix should not be read as a claim that one dataset is universally easier than another.
Each column uses the metric appropriate for that target dataset, and each source adapter may have learned a different output style.
For example, a ScienceQA rationale model may learn to produce reasoning before choosing an option, while a DocVQA answer-only model may learn short extractive text answering.
When these adapters are moved across domains, the result depends on both the learned objective and the target prompt.
This is why the transfer result is discussed as a pattern, not as a single ranking.

The most useful part of transfer evaluation is that it exposes when a gain is narrow.
A model can improve strongly in-domain but transfer poorly because it has adapted to the source dataset format.
Another model can have a smaller in-domain gain but transfer better to related visual formats.
For example, document and chart tasks can share some text-heavy or layout-heavy behaviour even though the images look different.
This kind of result helps prevent the thesis from treating a single benchmark as the whole story.

The transfer setting also explains why no-rationale datasets still matter.
ChartQA, DocVQA, and VQAv2 do not support rationale-supervised training in the completed setup, but they still test whether the same fine-tuning framework works outside rationale-heavy science and commonsense tasks.
They also give target domains for transfer analysis.
So those datasets are not unnecessary; they answer a different part of the thesis.
They show where answer-only adaptation and source-trained adapters work or fail when the visual domain changes.

\section{Reproducibility and artifact record}
Because the experiments are capped and single-seed, the artifact record is important.
The thesis cannot rely on the claim that results are obvious from large-scale repeated runs.
Instead, every completed run has to be traceable: which configuration was used, which dataset split was selected, which checkpoint was evaluated, which metric was computed, and which table or figure used the result.
This is why the pipeline writes an experiment ledger, a split audit, result tables, generated figures, and a claim-to-evidence map.

The experiment ledger records the backbone, dataset, objective, train and evaluation caps, prompt family, rationale availability, and primary metric.
This makes the comparison readable before the results are shown.
The split audit records the selected source split and checks whether selected training and evaluation identifiers overlap where native identifiers are available.
The purpose is to avoid accidental leakage and to make the capped protocol honest.
If a dataset source exposes only validation and test splits, the pipeline records exactly how the controlled split was built.

The claim-to-evidence map serves a different purpose.
It links each important claim in the abstract, highlights, discussion, and conclusion to a result row, table, figure, or artifact.
This is useful because it stops the paper from making a broader claim than the results support.
For example, a masking result can support a statement about evidence sensitivity, but it cannot by itself prove faithful explanation generation.
Likewise, a DocVQA answer-only result can support a claim about capped document answering, but it cannot support a claim about document rationale faithfulness.

The final validation step checks for missing expected runs, stale results, failed runs, and missing artifacts.
This does not make the experiments statistically stronger, but it makes the reporting cleaner.
It reduces the chance that a table is built from an old configuration or that a figure silently disappears after a partial run.
For a small thesis project, this kind of bookkeeping is part of research quality.
It keeps the scope modest, but it makes the completed scope easier to defend.
